%--------------------------------------------------------------------------
\chapter{Introduction (WORK)}
%--------------------------------------------------------------------------

%--------------------------------------------------------------------------
\section{Thermal Transport in Disordered Materials (WORK) (READ)}
%--------------------------------------------------------------------------

Why disordered materials are important...

Due to their potentially low thermal conductivities, 
disordered materials (e.g., alloys, amorphous solids, aerogels) 
are used in 
applications such a photoelectric energy conversion, 
thermoelectric energy conversion, and thermally insulating barriers.
\cite{graebner_phonon_1986,cahill_lattice_1988,lu_thermal_1992,
chen_recent_2003,clarke_thermal_2005,snyder_complex_2008,
minnich_bulk_2009,toberer_phonon_2011,zebarjadi_perspectives_2012,
schiffres_gas_2012,duda_exceptionally_2013}

Thermal management (i.e., heat spreaders and thermal insulators) 
depends critically on accurate predictions of thermal conductivity. 
Often the scattering from defects and boundaries reduces the thermal 
conductivity in these materials and devices, which is an area of 
active research.(cite) 

Recently, the search for more efficient thermoelectric 
materials has focused on reducing the thermal conductivity by alloying.
(cite)

The study of disordered lattices (i.e., alloys) has become increasingly 
quantitative as research seeks lower thermal conductivity 
thermoelectric materials, which have increased thermoelectric 
efficiency.(cite) Alloy modeling can be experimentally accurate 
to a high precision using quantum mechanically based density 
functional theory (DFT) calcualtions.(cite) However, DFT calcualtions 
are computationally expensive and are based on a unit cell 
with a small number of atoms.(cite) These calcualtions based on 
a small unit cell include the effects of disorder through 
perturbation techniques, which are only suited to a small amount of 
disorder. There is a need to asses the applicability of the perturbative 
methods used in DFT calculations so they can continue to make 
experimentally-accurate predictions. 

Amorphous materials have been long-studied by the physics community for 
their interesting properties such as the Boson peak...(cite)
However, most of thesis studies do not focus on thermal mode properties 
or conductivity predictions. 
Only in the past 20 years has the theory for the mode 
thermal properties been formulated.(cite) Most subsequent modeling 
studies have focused on amorphous silicon.



Thermal transport in ordered (crystalline) materials is well understood 
in terms of the phonon gas model.(cite) 
The focus of this work is on dielectric or semiconducting solids,  
where the heat is conducted by the atomic vibrational modes. 
Predicting the thermal conductivity of such materials 
requires the properties of the full spectrum of vibrational modes.
\cite{ziman_electrons_2001,feldman_thermal_1993,allen_diffusons_1999} 
Theoretical formulations for phonon properties date back to the work of 
Callaway,\cite{callaway_model_1959} 
Holland,
\cite{holland_analysis_1963} 
Klemens,\cite{klemens_scattering_1955} 
and Slack.\cite{slack_thermal_1979}. Their theory derived and/or 
postulated analytical models for the phonon dispersion and lifetimes 
(typically based on low-frequency limits) to be used in a solution 
of the BTE for predicting thermal conductivity. By fitting the
BTE solution to experimental bulk thermal conductivity data, expressions 
for the lifetimes are obtained. While good fits to the experimental 
data can be found, this agreement may be due to the numerous fitting 
parameters and not due to the correctness of the phonon lifetime models.
\cite{chung_role_2004} 
Atomistic techniques, which can predict the properties of individual phonon
modes without any assumptions about the dispersion or scattering, 
are thus required.

Newly-developed $\emph{ab initio}$ 
computational methods can make experimentally-accurate predcitions 
of the phonon properties and thermal conductivity of bulk crystals, 
whose structures are based on a unit cell with a small number of 
atoms.
\cite{broido_inrinsic_2007,ward_ab_initio_2009,intrinsic_ward_2010,
li_thermal_2012,lindsay_thermal_2012,lindsay_first-principles_2013}
The $\emph{ab initio}$ methods have been used to make predictions 
for materials which are 
difficult to produce experimentally, such as isotopically 
pure bulk,\cite{lindsay_thermal_2012} 
nanostructured\cite{kundu_light_2011,li_thermal_2012}
and difficult-to-manufacture 
materials.\cite{kundu_light_2011,lindsay_first-principles_2013} 
The lack of comparison with experiments 
in these $\emph{ab initio}$ 
studies makes it difficult to asses their predictive capability.
\cite{kundu_light_2011,lindsay_first-principles_2013} 

Incorporating 
disorder into the $\emph{ab initio}$ methods is difficult because 
a disordered material has a unit cell with an infinite 
number of atoms, by deinition. 
Disordering a crystal (i.e., alloying, amorphization) 
breaks down the phonon gas theory and a new theoretical formulation 
is necessary.\cite{allen_thermal_1993,allen_diffusons_1999} 
For small disorder, theoretical formulations begin 
with perturbation theory, which is valid for weakly disordered systems 
(i.e., isotopically disordered crystals).
\cite{klemens_scattering_1955,abeles_lattice_1963,tamura_isotope_1983} 
These perturbation theories can lead to simple models which can be 
fit empirically to experimental measurements. However, predictive 
methods which do not rely on perturbation theory are still in active  
development.  
Comparison with experiment is excellent for recent 
DFT predictions of thermal conductivity. 
For disordered materials, the comparison to experiment is often 
only qualititative. 
\cite{shiomi_thermal_2011,tian_phonon_2012,tian_enhancing_2012}

Accurate modeling of amorphous materials is even more 
difficult with DFT calcualtions because the strong topological 
disorder prohibits the use of perturbation theory.(cite)  
A finite unit cell can be used to model an amorphous material, but 
the size of these modes are typically larger than 100 atoms.(cite) 
The theoretical framework for studying heat transport in 
amorphous materials is fundamentally different from 
phonon perturbation theory.
\cite{allen_thermal_1993,allen_diffusons_1999} No comprehensive 
study has been performed to asses the applicability of the 
phonon perturbation and amorphous thermal transport theories 
on materials with a range of disorder (i.e., crystals, 
weakly disordered crystals, heavily disordered crystals, 
and amorphous materials). 

Thermal transport in disordered materials is studied in this 
work using several predictive methods (see Section \ref{ThermalModes}). 
Because disordering complicates 
the theory of thermal transport, no single method can provide 
both the system-level (total thermal conductivity) and mode-level 
properties for all vibrations in the system. Each 
predictive method provides complimentary information which is 
compared wherever possible. Important 
assumptions about thermal transport in disordered 
materials are investigated using fully atomisitic models. 
Using these comparisons, critical information 
about the vibrational mode-level properties is identified.  
The predicted 
results are compared to phenomological and empirical 
models. The results presented in this work provide a theoretical and 
computational framework for the study of emerging disordered 
and nanoscaled systems, which are discussed in Section 
\ref{Conclusion}. 

%--------------------------------------------------------------------------
\section{\label{ThermalModes}
Thermal Modes: Phonons, Propagons, Diffusons, and Locons 
(WORK) (READ)}
%--------------------------------------------------------------------------

For a perfect lattice, 
all vibrational modes are phonon modes, which by 
definition are delocalized, propagating plane waves.
\cite{ziman_electrons_2001} For disordered materials, all vibrational 
modes are not phonons because there is no translational order.  
Allen and Feldman define all disordered vibrational modes to be 
vibrons because they are all bosons and follow Bose-Einstein 
statistics.\cite{allen_diffusons_1999} The vibrons are propagons, 
diffusons, and locons. Propagon are propagating and delocalized 
(i.e., phonon-like) and are typically found in the low-frequency 
range of the vibrational spectrum.(cite) 
Diffusons are non-propagating and delocalized modes which couple 
harmonically due to the disorder and spatial delocalization.
\cite{allen_thermal_1993} While diffusons are non-propagating, 
they can contribute a significant amount to thermal transport,
\cite{feldman_thermal_1993,feldman_numerical_1999}. 
Locons are non-propagating and spatially 
localized modes which do not couple to others, and are typically 
found at the high frequency of the vibrational spectrum. 
Because of this locons are spatially localized, 
they do not contribute to thermal transport.(cite) 

The spectrum of vibrons is depcited in Fig. 
\ref{F:allen_diffuson_1999_fig1} for a model of amorphous silicon.
\cite{allen_diffusons_1999} The propagons exist at low frequencies below 
the Ioffe-Regel limit,(cite) which marks the transition between 
propagating and non-propagating modes. The IR limit is discussed in 
Sections and . 
Propagaons and diffusons are called extendons because they are 
delocalized. The mobility edge marks the transition from 
the delocalized diffusons to the localized locons.(cite) 


%--------------------------------------------------------------------------
\begin{figure}
\begin{center}
\centering
\includegraphics[scale=0.4]
{/home/jason/thesis/thesis/allen_diffuson_1999_fig1.eps}
\end{center}
\caption{\label{F:allen_diffuson_1999_fig1} 
Density of vibrational states of amorphous silicon and classification 
of vibrational modes (vibrons). Figure is reproduced from 
Ref. \citenum{allen_diffusons_1999}.
}
\end{figure}
%--------------------------------------------------------------------------
\clearpage

%--------------------------------------------------------------------------
\section{\label{S-SystemMode}
System-level versus Mode-level Thermal Predictions}
%--------------------------------------------------------------------------

%--------------------------------------------------------------------------
\subsection{\label{S-ktopdown} 
System-level Thermal Conductivity Prediction (WORK) (READ)}
%--------------------------------------------------------------------------

% There are two commonly-used top-down techniques for predicting thermal 
% conductivity from MD simulations: the Green-Kubo method 
% \cite{mcquarrie2000,mcgaughey2006c}, which uses equilibrium simulations, 
% and the direct method \cite{schelling2002,mcgaughey2006c}, which uses 
% steady-state non-equilibrium simulations. These two methods are compared 
% in Table \ref{T-ktopdown}. While so-called ``quantum corrections" to 
% thermal conductivities predicted from classical MD simulations have been 
% proposed, Turney et al. demonstrate that these corrections are not 
% rigorous and should not be applied \cite{turney2009b}. Comparison of 
% MD-predicted thermal conductivities to experimental measurements should 
% therefore be limited to high temperatures, around and above a material's 
% Debye temperature. Top-down techniques are important in the context of 
% this chapter because they can be used to validate the bottom-up 
% techniques that are our focus.
% 
% The Green-Kubo and direct methods can be applied to crystals, alloys, 
% amorphous solids, and fluids. Recent work has focused on computational 
% aspects, notably accounting for size effects, quantifying and reducing 
% uncertainty, and reducing computational time 
% \cite{landry2008,zhou2009,jchen2010,
% sellan2010a,howell2011a,howell2011b,
% he2012,howell2012,jones2012}.
% The Green-Kubo method relates the equilibrium fluctuations
% of the heat current vector, $\mathbf{S}$, to the thermal conductivity 
% tensor
% via the fluctuation-dissipation theorem. The thermal conductivity in 
% direction $\alpha$ is given by \cite{mcquarrie2000}
% \begin{eqnarray}\label{E-GK}
% k_{\alpha} = \f{1}{k_\mathrm{B} V T^2} 
% \int_0^{\infty} \lr{S_{\alpha}(t)S_{\alpha}(0)}dt, 
% \end{eqnarray}
% where $k_{\mathrm B}$ is the Boltzmann constant, $t$ is time, and 
% $S_{\alpha}$ and $\lr{S_{\alpha}(t)S_{\alpha}(0)}$ are the $\alpha$ 
% components of the heat current vector and the heat current autocorrelation 
% function (HCACF). The heat current vector is calculated using the atomic 
% positions and velocities. When using the Green-Kubo method, we recommend 
% that the MD simulations be run in the $NVE$ ensemble.\footnote{As mentioned 
% in Section \ref{S-MDformulation}, $NVE$ simulations are advantageous 
% because the equations of motion are not modified. The Green-Kubo method 
% formula for thermal conductivity, Eq.~\eqref{E-GK}, however, is derived 
% in the $NVT$ ensemble, leading to the need to specify the temperature of 
% the MD simulation. Temperature is not well-defined in the $NVE$ ensemble, 
% but can be approximated using the system kinetic energy and an assumption 
% of equipartition of energy (very good for classical MD simulation of 
% solids).} The Green-Kubo method is advantageous for anisotropic materials 
% because it can be used to predict the full thermal conductivity tensor.

There are several atomistic methods for predicting the thermal 
conductivity at the system level. The molecular dynamics (MD)-based 
\href{http://lammps.sandia.gov/doc/fix_thermal_conductivity.html}
{direct method} predicts the thermal conductivity by a direct 
application of the 
\href{http://en.wikipedia.org/wiki/Thermal_conduction#Fourier.27s_law}
{Fourier law} for heat conduction.(cite) While the direct method 
is conceptually simple, it can be difficult to account for finite-size 
effects.(cite) The 
\href{http://lammps.sandia.gov/doc/compute_heat_flux.html}{Green-Kubo}
method can predict the thermal conductivity from equilibrium MD 
simulations. These equilibrium simulations are often simpler 
to setup and account for finite-size effects (see Section ).
(cite) The GK method has been used extensively in other modeling 
studies(cite) and is discussed in detail elsewhere.(cite) 
While the GK  and direct methods have been shown to take into 
account the effects of disorder,(cite)they predict no mode-level 
properties of the thermal carriers. 

The Green-Kubo and direct methods can be applied to crystals, alloys, 
amorphous solids, and fluids. Recent work has focused on computational 
aspects, notably accounting for size effects, quantifying and reducing 
uncertainty, and reducing computational time 
\cite{landry_complex_2008,buongiorno_benchmark_2009,
chen_how_2010,sellan_size_2010,he_lattice_2012}.
need howell2011a howell2011b howell2012 jones2012

While so-called ``quantum corrections" to 
thermal conductivities predicted from classical MD simulations have been 
proposed, Turney et al. demonstrate that these corrections are not 
rigorous and should not be applied \cite{turney_assessing_2009}. 
Comparison of MD-predicted thermal conductivities to experimental 
measurements should 
therefore be limited to high temperatures, around and above a material's 
Debye temperature. The GK method, which is a system-level technique, 
is important for the present study because it can be used to validate 
the mode-level techniques that are the focus. The capabilities 
of the GK method are summarized in Table \ref{T-comparison}.


%--------------------------------------------------------------------------
\subsection{\label{S-kbottomup} Mode-level Thermal Conductivity 
Prediction (WORK) (READ)}
%--------------------------------------------------------------------------

While the total thermal conductivity is typically the quantity of 
interest for modeling and engineering analysis, the mode-level 
properties can be just as important. The mode-level properties 
allow for a deeper understanding of the total thermal conductivity.
(cite) 
The mode-level properties are used to understand the difference 
between predictions for alloys in Chapter and amorphous 
materials in Chapter . The mode-level properties also allow for an 
understanding of the effects of nanostructuring, which can cause 
increased scattering from boundaries(cite) and modifications to the 
vibrational density of states.(cite)

The Callaway-Holland (CH) theory, which is based on the work of 
Callaway,\cite{callaway_model_1959} Holland,
\cite{holland_analysis_1963} Klemens,\cite{klemens_scattering_1955} 
and Slack,\cite{slack_thermal_1979} uses analytical models for phonon 
dispersion and lifetimes. These analytical models are often empircally 
fit to experimental results, which limits the predictive capability of 
the CH theory. While the CH theory can calculate mode properties 
from the analytical models, it is unclear how accurate these properties 
are.

There are several methods which can predict the mode-level vibrational 
properties. 
Anharmonic lattice dynamics (ALD) calculations,
\cite{turney_predicting_2009,esfarjani_heat_2011} 
based on perturbation 
theory,\cite{maradudin_dynamical_1974,srivastava_physics_1990} 
can be used along with phonon-defect scattering perturbation 
theory to predict the mode properties and thermal conductivity 
of mass and/or bond disordered lattices (see Section ).
\cite{klemens_scattering_1955,klemens_thermal_1957,mattis_phonon_1957,
tamura_isotope_1983} 
However, the ALD is a perturbative method which is valid at low 
temperatures where 3-phonon processes dominate.(cite) 
Because the ALD method is based on a unit cell 
and disorder is included as a perturbation, the mode properties 
predicted are for phonons. One of the objectives of this work is 
to investigate if the phonon properties predicted by the ALD method 
are consistent with the properties of propagaons in a fully 
disordered system (see Section ). 

The ALD and phonon-defect perturbative method is computationally 
inexpensive so that 
computationally expensive DFT calcualtions can be used to predict 
thermal conductivities that are experimentally accurate.
\cite{garg_role_2011,lindsay_thermal_2012,li_thermal_2012} 
However, most studies using ALD and DFT have been limited to 
disordered lattices (i.e., alloys) whose thermal conductivities 
are large and dominated by low-frequency propagons.
(cite) 
Because no comprehensive study has been performed on a range 
of materials, it is unclear what limitations exist for using the ALD and 
phonon-defect perturbative method.
\cite{luckyanova_coherent_2012,tian_phonon_2012,tian_enhancing_2012} 
We investigate the limitations of the ALD and phonon-defect perturbative 
method in Section . 

The key to explicitly incorporating the effects of disorder 
is to use large disordered supercells, or a unit cell with a large 
number of atoms, $N_a$.  
The use of computationally 
inexpensive empirical potentials allows for calcualtions to be 
performed on a large number of atoms, typically $N_a$ > 500.(cite) 
The Allen-Feldman (AF) theory computes 
the contribution of diffusons 
to thermal conductivity from large disordered supercells.
\cite{allen_thermal_1993} 
The AF theory calculations are based on harmonic lattice dynamics which 
rely on the system disorder to couple modes thermally.
\cite{allen_thermal_1993,feldman_thermal_1993,feldman_numerical_1999} 
However, because the AF theory is harmonic, it is not suited to studying 
the low-frequency propagons in disordered materials. For disordered lattices, 
the low-frequency propagons are subject to Rayleigh scattering under the 
harmonic approximation, which 
causes a divergence in the thermal conductivity in the zero frequency 
limit.(cite)  
The Allen-Feldman (AF) theory can use quantum or classical expressions 
for the specific heat. While the AF theory should be limited to low 
temperatures where the harmonic approximation is valid, 
anharmonic effects have been shown to be 
small for disordered solids such as amorphous silicon.
\cite{feldman_thermal_1993}

The NMD method can be used to predict the lifetime of all vibrons from 
MD simulations.
\cite{bickham_calculation_1998,
bickham_numerical_1999,fabian_numerical_2003,
ladd_lattice_1986,mcgaughey_quantitative_2004,donadio_thermal_2007,
henry_spectral_2008,donadio_atomistic_2009,turney_predicting_2009,
goicochea_thermal_2010,he_morphology_2011,he_thermal_2011,he_heat_2011,
he_lattice_2012,sosso_thermal_2012,qiu_molecular_2012,wang_mode-wise_2013,
bao_first-principles_2013} 
While the method can only predict the 
mode lifetimes, these can be used with other predictive methods 
and assumptions to predict the thermal conductivity (see Sections ). 
NMD is limited to mid-range temperatures because of the harmoninc 
approximation used to calcualte the normal mode eigenvectors. 
The thermal conductivity predicted 
by NMD begins to deviate from GK and direct-method prediction 
near about half the melting temperature for lennard-jones argon.
\cite{turney_predicting_2009} 

Because of the theroetical limtations of the methods discussed above, 
no one method is sufficient for predicting the mode 
properties of all vibrons in a given disordeded material. The 
theoretical methods are summarized and compared in 
Table \ref{T-comparison}. 
The methods in Table \ref{T-comparison} are ranked in the order 
of their thermal conductivity predictive capability based on 
prior work in the literature.(cite) 
For example, the CH theory relies completely on empircal fitting, 
while the GK method can predict the thermal conductivity of some 
materials in any phase (crystalline, amorphous, liquid).(cite) 
In Section , the predictive capabilities are ranked in 
Table \ref{T-comparison-predictive} based on the results from 
this work.

(more discussion on Table \ref{T-comparison})

%--------------------------------------------------------------------------
\begin{center}
%\squeezetable
\begin{table}
\small
\caption{\label{T-comparison}Comparison of theoretical 
techniques for predicting vibrational mode properties and thermal 
conductivity for disordered systems.}
%\begin{tabular}{llllll}
\begin{tabular}{p{0.9in}|p{0.9in}|p{0.9in}|p{0.9in}|p{0.9in}|p{0.9in}}
\hline\hline
&CH Theory 
&ALD
&AF Theory
&NMD
&GK \\ 
\hline\hline
Description
&Empirical
&Predictive 
&Predictive 
&Predictive
&Predictive\\
\hline
Input
&Fitted Models
&Force constants \newline (classical or DFT)  
&Force constants \newline (classical or DFT)
&Forces and Force Constants \newline (classical or DFT) 
&Forces \newline (classical or DFT) \\
\hline
Disorder
&As additional \newline perturbation 
&As additional \newline perturbation 
&Naturally included 
&Naturally included 
&Naturally included \\
\hline
Statistics
&Bose-Einstein or \newline Boltzmann 
&Bose-Einstein or \newline Boltzmann 
&Bose-Einstein or \newline Boltzmann
&Boltzmann 
&Boltzmann\\
\hline
Temperature
&All 
&Low 
&All
&Mid-range
&All\\
\hline
Modes
&Phonons (Propagons?)
&Phonons (Propagons?)
&Diffusons
&Propagaons and Diffusons
&N/A\\
\hline \hline
\end{tabular}
\end{table}
\end{center}
%--------------------------------------------------------------------------
\clearpage

%--------------------------------------------------------------------------
\subsection{\label{S-CompCost}
Computational Cost vesus Preditictive Ability (WORK) (READ)}
%--------------------------------------------------------------------------

The computational costs of the methods discussed in Sections 
\ref{S-ktopdown} and \ref{S-kbottomup}  
are given in Table \ref{T-comparison-cost}.  
In most cases, the 
computational cost of a method is known $\emph{a priori}$ and should 
be considered before any analysis is performed. The computational 
costs are for these methods can be represented by the total number of 
atoms, $N_a$, that are being simulated. For ordered and disordered 
systems, the total 
number of atoms can be related to the number of atoms in the unit 
cell ($n$) and the 
number of unit cells, $N$ ($N$ = 0 for a disordered system), 
which gives $N_a = N n$ (see Section ).

One of the goals 
of this work is to more thoroughly asses the predictive 
abilities of these methods.
The predictive abilities of all methods are examined throughout 
this work and are summarized in Section \ref{S:Predictive}, 
Table \ref{T-comparison-predictive}. 
The NMD method can predict the mode lifetimes for the vibrational 
modes in ordered and disordered systems, making its predictive 
ability high.
\cite{ladd_lattice_1986,mcgaughey_quantitative_2004,donadio_thermal_2007,
henry_spectral_2008,donadio_atomistic_2009,turney_predicting_2009,
goicochea_thermal_2010,he_morphology_2011,he_thermal_2011,he_heat_2011,
he_lattice_2012,sosso_thermal_2012,qiu_molecular_2012,wang_mode-wise_2013,
bao_first-principles_2013} 
However, the explicit use of disorder makes it 
computational cost high (see Section \ref{Subsection_Comp_Details_4}). 
Efficient MD 
codes like LAMMPS scale linearly with the number of atoms in 
the system, $N_a$, which makes the GK method (see Section 
\ref{S:Thermal Conductivity}) 
computationally-inexpensive when used with classical interatomic 
potentials.(cite) 
 
For ALD, $N_a$ determines the resolution of the 
\href{https://en.wikipedia.org/wiki/Brillouin_zone}
{Brillouin zone} sampling (see Section \ref{Subsection_Comp_Details_1} 
and Ref. \citenum{turney_predicting_2009-1}). Some systems (such as 
the silicon alloys studied in Section \ref{S:SW}) 
require a fine sampling 
of the BZ, which increases the ALD computational costs. Using large 
systems can be avoided by following the extrapolation procedure 
which is discussed in Appendix \ref{Appendix_A:Finite}.

The ranking of computational cost in Table \ref{T-comparison-cost} 
does not take into account the coefficients of the computaional 
scalings. For example, to perform NMD 
on a system of 6912 atoms in less than 24 hours may take 100 
CPUs (see Section ). The same number of CPUs are required to perform 
the GK method on a system of 100000 atoms in the same amount of time. 
Coefficients of the computational scalings would take into 
account the difference in computational cost of classical pair 
potentials versus multi-body potentials, for example. 
Thus, the ranking in Table \ref{T-comparison-cost} is to be used as a 
rough guideline. 

(more discussion of the Table \ref{T-comparison-cost})

The GK method, which scales most efficiently in both memory and CPU 
cost, has the most potential to simulate systems with increased 
$N_a$. While $N_a$ are reached with the ALD method in this work, 
no explicit atomic dynamics are simulated. 

%--------------------------------------------------------------------------
\begin{center}
%\squeezetable
\begin{table}
\small
\caption{\label{T-comparison-cost}Ranking of computational costs 
for theoretical 
techniques for predicting vibrational mode properties and thermal 
conductivity.}
%\begin{tabular}{llllll}
\begin{tabular}{p{0.9in}|p{0.9in}|p{0.9in}|p{0.9in}|p{0.9in}|p{0.9in}}
\hline\hline
&CH Theory 
&AF Theory
&NMD
&GK
&ALD\\ 
\hline
Computational Cost\\
\hline
CPU\newline
Memory 
&Minimal \newline Minimal
&High \newline High 
&High \newline High 
&High \newline Low 
&High \newline Low \\
\hline
Maximum System Size
&N/A
&6912 atoms
&6912 atoms
&592704 atoms
&800000 atoms\\
\hline
Computational Scaling\\
\hline
CPU\newline
Memory 
&Minimal \newline Minimal
&$(Nn)^{3}$ \newline $(Nn)^{2}$ 
&$(Nn)^{2}$ \newline $(Nn)^{2}$ 
&$(Nn)^{1}$ \newline $(Nn)^{1}$ 
&$N^{2}n^{4}$ \newline $Nn^{2}$  \\
\hline\hline
\end{tabular}
\end{table}
\end{center}
%--------------------------------------------------------------------------
\clearpage